\pg{359} \emph{mutavel é possivel}, ou está em potencia em relação ao acto, para o qual se muda. Logo o. ente \emph{mutavel} deve ser reduzido ao acto por outro ente. --- Este outro ente ou é immutavel ou é mutavel: Se éimmnutavel; está provada a conclusão. Se é mutavel, é \emph{possivel} e por isso suppõe um outro ente, e este um outro, até se encontrar um immutavel; pois, se não existisse o ente immutavel, não poderia existir nenhum ente mutavel. Logo o ente \emph{mutavel} suppõe o ente immutavel (».

90. Corollarios. --- Dos principios expostos resulta :

T. Na ordem \emph{ontologica}, ou da \emph{realidade}, o immutavel precede o \emph{muíavel}; mas, na ordem \emph{logica}, ou do \emph{conhecimento}, o \emph{mutavel} precede o \emph{immutavel}, porque é do conhecimento d'aquelle que nos elevamos ao conhecimento d' este. m

IT. Dada a existencia de um só ente \emph{mutavel}, deve neces- riamente admittir-se a existencia de um ente immutavel \footnote{Vemos no mundo entes sujeitos a mudanças; que do não-ser passa- \textasciicircum{} \textasciicircum{} ram para o ser. Esta passarem não poude ter sido realisada pelos proprios entes, que. estando em potencia, isto é, sendo possiveis mas não extistentes, não podiam operar; pois \& operação é uma consequencia do ser. Logo essa passagem deve ter a sua causa proporcionada na força de um ente exterior , e, em ultima analyse. absolutamente immutavel: porque. se não existisse o \emph{imnmutavel}, não poderia existir. como effectivamente existe, o mutavel. Diz S. Thomaz: « Videmus aliquid esse in mundo. quod exit de potentia in aclum. Non autem educit se de potentia in actum : quia quod est (in) potentia. nondum est, unde nec agere potest. Ergo oportet esse aliquid aliud Prius, quo educatur de potentia in actum. Et iterum, si hoc est exiens de potentia in actum, oportet ante hoc aliquid aliud poni, quo reducatur in . actum. Hoc autem jn infinitum procedere non polest. Ergo oportet devenire ad aliquid, quod est tantum actu, et nullo modo in potentia; et hoc dicimus \textsc{Deum}» (\emph{C}. \emph{Gent}., I. 16). .}...\footnote{Diz S. Fulgencio: «Fjrmissime tene et nullatenus dubites omnem creaturam naturaliter mutabilem a \textsc{Deo} incommutabili factam » (\emph{De fide ad Petrum}, c. 90). --- Recommendamos aos leitores o estudo da doutrina, exposta neste capitulo segundo ; pois constitue ella a base dos argumentos, que provam a existencia de Deus. !}
